% Source-aligned English for printed p.166.
\SourcePage{166}

\LayerTarget{tohoku:remark:3.4.2:3:p2}
structure sheaf of $E$ to reduce to the \textit{constant} sheaf $G$, that is,
for $E$ to admit an integrable holomorphic connection. This condition is
sufficient when $X$ has complex dimension 1, for then $u$ is surjective. The
invariant $d(E)$ was first introduced by Weil [21]; a more geometric
definition due to Atiyah [1] shows that the vanishing of $d(E)$ is in every
case necessary and sufficient for the existence on $E$ of a holomorphic
connection, not necessarily integrable.

\LayerTarget{tohoku:section:3.5}
\indent\textbf{3.5. The exact sequence relative to a closed subspace.} Let $Y$ be a
locally closed subspace of $X$, that is, the intersection of an open subset
and a closed subset. For every abelian sheaf $F$ on $Y$, there exists a unique
abelian sheaf on $X$ whose restriction to $Y$ is $F$ and whose restriction to
$\complement Y$ is zero. To see this, one reduces at once to the case in which $Y$ is
open or closed; the straightforward verification is given in [4, XVII,
Prop. 1].
\LayerTarget{tohoku:def:extension-by-zero-locally-closed}
Denote this sheaf on $X$ by $F^X$. The assignment
$F\to F^X$ is an exact functor $\mathbf C^Y\to \mathbf C^X$. Moreover,
if $Z\subset Y\subset X$, with $Z$ locally closed in $Y$ and hence in $X$,
and if $F$ is an abelian sheaf on $Z$, then
\LayerTarget{tohoku:eq:extension-by-zero-transitivity}
$(F^Y)^X=F^X$.
\LayerTarget{tohoku:def:restriction-extension-FY}
Now let $F$ be an abelian sheaf on $X$ and put $F_Y=(F\vert Y)^X$. Thus $F_Y$
is the sheaf on $X$ characterized by having the same restriction to $Y$ as
$F$ and zero restriction to $\complement Y$. The assignment $F\to F_Y$ is an
exact functor $\mathbf C^X\to \mathbf C^X$. It also has the transitivity property
\LayerTarget{tohoku:eq:restriction-extension-transitivity}
$(F_Y)_Z=F_Z$
when $Z\subset Y\subset X$ are as above. If $Y$ is closed, so that $U=\complement Y$ is
open, there is the familiar exact sequence

\LayerTarget{tohoku:eq:closed-subspace-exact-sequence}
\[
0\to F_U\to F\to F_Y\to0
\]
for every abelian sheaf $F$ on $X$. Recall that
$\Gamma(F_Y)=\Gamma(Y,F)=\Gamma(F\vert Y)$, whereas $\Gamma(F_U)$ identifies
with the subgroup of $\Gamma(F)$ formed by sections whose support is contained
in $U$.
\LayerTarget{tohoku:def:induced-support-antifilter}
Let $\Phi$ be an antifilter of closed subsets of $X$. For every subset $Z$ of
$X$, let $\Phi_Z$ be the ``induced'' antifilter formed by the $A\in\Phi$ that
are contained in $Z$; this is not to be confused with the trace
$\Phi\cap Z$ of $\Phi$ on $Z$. If $Z$ is locally closed and $\Phi$ is
paracompactifying, as in 3.3, then $\Phi_Z$ is also paracompactifying. If $Z$
is locally closed and $\Phi$ is arbitrary, the formulas above give the more
general identity below, valid in particular when $Z$ is open or closed:

\LayerTarget{tohoku:eq:sections-support-restriction}
\[
\Gamma_\Phi(X,F_Z)=\Gamma_{\Phi_Z}(Z,F\vert Z)
\]
for every abelian sheaf $F$ on $X$. Equivalently,

\LayerTarget{tohoku:eq:sections-extension-by-zero}
\[
\Gamma_{\Phi_Z}(Z,G)=\Gamma_\Phi(X,G^X)
\]
for every abelian sheaf $G$ on $Z$.
\LayerTarget{tohoku:construction:cohomology-support-comparison:p1}
Since $G\to G^X$ is \textit{exact}, the groups $H^p_\Phi(X,G^X)$ form a
cohomological functor on $\mathbf C^Z$. The universal cohomological functor
$(H^p_{\Phi_Z}(Z,G))$ agrees with it in degree 0, and therefore gives
\textit{canonical homomorphisms}

\LayerTarget{tohoku:eq:canonical-support-comparison}
\[
H^p_{\Phi_Z}(Z,G)\to H^p_\Phi(X,G^X)
\]
(characterized by defining a homomorphism of cohomological functors

\TohokuContinuation{167}
